\begin{lemma}[Strong law of large numbers for a Paolini--Stepanov family, single stage]\label{SLLNCurves}
  Let $E$ be a metric space equipped with a compatible Borel $\sigma$-algebra,
  $T \colon \mathcal{D}^1(E) \to \mathbb{R}$, and $F$ a Paolini--Stepanov family for $T$
  (\noderef{PSFamily}). Fix a stage $l \in \mathbb{N}$, $l \ge 1$, with
  $\overline{\eta}_l(\Theta(E)) \ne 0$ (\noderef{PSFamily}, measure $\overline{\eta}_l$).
  Let $(\omega_k)_{k \in \mathbb{N}}$ be an arbitrary countable family of metric $1$-forms
  (\noderef{Form1}), and let $(g_j)_{j \in \mathbb{N}}$ be an arbitrary countable family of
  measurable functions $g_j \colon \Theta(E) \to \mathbb{R}$, each $\overline{\eta}_l$-integrable.
  Then there is a single sequence of curves $(\gamma_i)_{i \in \mathbb{N}} \subseteq \Theta(E)$
  such that:
  \begin{itemize}
    \item[(i)] $\length(\gamma_i) = l$ for every $i$;
    \item[(ii)] $\mathbb{M}([[\gamma_i]]) = l$ for every $i$ (equivalently $\gamma_i \in
      \Theta_l(E)$, \noderef{ThetaL});
    \item[(iii)] for every $k \in \mathbb{N}$,
      \[
        \frac{\overline{\eta}_l(\Theta(E))}{n \cdot l} \sum_{i=0}^{n-1} [[\gamma_i]](\omega_k)
        \;\longrightarrow\; T(\omega_k)
        \qquad (n \to \infty)
        ;
      \]
    \item[(iv)] for every $j \in \mathbb{N}$,
      \[
        \frac{1}{n} \sum_{i=0}^{n-1} g_j(\gamma_i)
        \;\longrightarrow\; \overline{\eta}_l(\Theta(E))^{-1} \int_{\Theta(E)} g_j \, d\overline{\eta}_l
        \qquad (n \to \infty)
        .
      \]
  \end{itemize}
  Here $\overline{\eta}_l(\Theta(E))$ (a finite, nonzero real number, since $\overline{\eta}_l$ is a
  finite measure by hypothesis) denotes the real number obtained from the extended-real
  $\overline{\eta}_l(\Theta(E))$ by \texttt{ENNReal.toReal}, and the sums in (iii)--(iv) are
  indexed $i=0,\dots,n-1$ (\texttt{Finset.range n}); this is a pure reindexing of the paper's
  $i=1,\dots,n$ convention (paper.tex 1295--1310) and carries no mathematical difference.

  \paragraph{Role of this lemma.} This is the Strong Law of Large Numbers step of the proof of
  Theorem~4.4 (\eqref{countable_identity} and \eqref{stronglawlength}, paper.tex 1295--1310),
  restated at a single fixed stage $l$ and for an arbitrary pair of countable families
  $(\omega_k)_k$, $(g_j)_j$ (rather than the paper's specific families $\tilde f_k, \tilde\pi_j$
  and the marginal-cutoff / $\psi$-average functions), so that \noderef{BBSamplingDiagonal} may
  instantiate $(g_j)_j$ with whichever specific $\overline\eta_l$-integrable countable family it
  needs (the endpoint-cutoff indicators and the $\psi$/$\tilde\psi$-averages) without this node
  itself naming them. Clause (iii) is the substantive content beyond
  \noderef{SamplingRealizesAverages}: it exhibits the explicit
  $\overline{\eta}_l(\Theta(E))$ (later identified with $\mathbb{M}(T)$, an identification this
  node does not need and does not make) that eq.~(4.6) forces into the normalizing constant of
  \eqref{closed1}, and does so by combining \noderef{SamplingRealizesAverages}, applied once, with
  \noderef{PSFamily}'s \emph{weakLength} field. Clause (iv) is the generic (unnormalized-except-by
  $\overline\eta_l(\Theta(E))$) strong-law average, stated once so that every other countable
  strong-law family the assembly needs is obtained by instantiating $g$, rather than by re-running
  \noderef{SamplingRealizesAverages} again.

  \paragraph{Design point: a single, internally-merged application of
  \noderef{SamplingRealizesAverages}.} Both the current-family $(\omega_k)_k$ and the auxiliary
  family $(g_j)_j$ are folded into one $\mathbb{N}$-indexed family $h$ (even indices carry the
  currents against $\omega_k$, odd indices carry $g_j$) \emph{inside the proof}, so that
  \noderef{SamplingRealizesAverages} is invoked exactly once, on $h$, rather than requiring a
  separate packaging node for the merge. Integrability of $\gamma \mapsto [[\gamma]](\omega_k)$
  with respect to $\overline\eta_l$ is likewise established inside this proof (from
  \noderef{CurveCurrentBound} together with \noderef{PSFamily}'s \emph{lengthExactly} field and
  finiteness of $\overline\eta_l$), so it is not a hypothesis of the statement.

  \paragraph{Design point: $\Theta_l(E)$ need not be shown measurable.} The set $\Theta_l(E)$
  (\noderef{ThetaL}) is a condition on the mass of a curve's current and is not known to be
  measurable with respect to the evaluation $\sigma$-algebra on $\Theta(E)$ this tablet uses, nor
  is that needed here: \noderef{PSFamily}'s \emph{supported} field only gives that the
  \emph{complement} of $\Theta_l(E)$ is $\overline\eta_l$-null, and Mathlib's
  \texttt{MeasureTheory.exists\_measurable\_superset\_of\_null} produces, from any
  $\overline\eta_l$-null set, a \emph{measurable} superset $U$ that is still $\overline\eta_l$-null;
  taking $S := U^c \cap \{\gamma : \length(\gamma) = l\}$ (measurable, since $U$ is measurable and
  $\{\length(\cdot)=l\}$ is measurable by \noderef{CurveLenMeasurable}) gives a measurable,
  co-null, subset of $\Theta_l(E) \cap \{\length(\cdot)=l\}$ to feed
  \noderef{SamplingRealizesAverages} without ever asking whether $\Theta_l(E)$ itself is
  measurable.

  \paragraph{Design point: the only nondegeneracy hypothesis is $\overline\eta_l(\Theta(E)) \ne
  0$.} No hypothesis of \noderef{PaoliniStepanovExists} or \noderef{MassPosOfNonzero} is needed at
  this stage; the downstream target-covering assembly is responsible for discharging
  $\overline\eta_l(\Theta(E)) \ne 0$ later, from \noderef{PSFamily}'s \emph{massMarginal} field at
  $\phi \equiv 1$ together with \noderef{PaoliniStepanovExists} and \noderef{MassPosOfNonzero}. The
  identity $\overline\eta_l(\Theta(E)) = \mathbb{M}(T)$ (which is where eq.~(4.6)'s implicit
  $\mathbb{M}(T)$ factor of \eqref{closed1} ultimately comes from) is likewise established
  downstream, not here.

  \paragraph{Why no separate $[\mathrm{Nonempty}\ E]$ hypothesis is needed.} The conclusion
  asserts the existence of curves $\gamma_i \in \Theta(E)$, and $\Theta(E)$ is empty exactly when
  $E$ is (\noderef{Curve} carries $\mathrm{toFun}\colon\mathbb{R}\to E$, and no total function from
  the nonempty domain $\mathbb{R}$ into an empty codomain exists). But if $E=\emptyset$ then
  $\Theta(E)$ itself is empty, so $\mathrm{Set.univ}\subseteq\Theta(E)$ is the empty set, and every
  measure assigns the empty set mass $0$; hence $\overline\eta_l(\Theta(E)) =
  \overline\eta_l(\emptyset) = 0$ automatically whenever $E=\emptyset$, contradicting the standing
  hypothesis $\overline\eta_l(\Theta(E))\ne0$. So $E=\emptyset$ is already excluded, and no
  additional $[\mathrm{Nonempty}\ E]$ is required.
\end{lemma}
\begin{proof}
  Write $M := \overline\eta_l(\Theta(E)) \in \cointerval{0,\infty}$ for the total
  $\overline\eta_l$-mass. By hypothesis $M \ne 0$, and $M \ne \infty$ since $\overline\eta_l$ is a
  finite measure (\noderef{PSFamily}, field \emph{finite}); write $M_{\mathbb R} :=
  \mathrm{toReal}(M) \in (0,\infty)$ for its real value, so $M_{\mathbb R} \ne 0$ and
  $M^{-1} \ne \infty$.

  \emph{Step 1: integrability of the current family.} Fix $k \in \mathbb{N}$. By
  \noderef{CurveCurrentMeasurable}, $\gamma \mapsto [[\gamma]](\omega_k)$ is measurable. By
  \noderef{PSFamily} (field \emph{lengthExactly}), $\length(\gamma) = l$ for $\overline\eta_l$-a.e.\
  $\gamma$; on this event, \noderef{CurveCurrentBound} gives
  $\bigl|[[\gamma]](\omega_k)\bigr| \le \mathrm{fBound}(\omega_k)\cdot\mathrm{piLip}(\omega_k)\cdot
  \length(\gamma) = \mathrm{fBound}(\omega_k)\cdot\mathrm{piLip}(\omega_k)\cdot l$, a constant.
  A measurable function that is $\overline\eta_l$-a.e.\ dominated by an $\overline\eta_l$-integrable
  constant function (constants are integrable against a finite measure) is itself
  $\overline\eta_l$-integrable; hence $\gamma \mapsto [[\gamma]](\omega_k)$ is
  $\overline\eta_l$-integrable, for every $k$.

  \emph{Step 2: the merged family.} Define $h \colon \mathbb{N} \to \Theta(E) \to \mathbb{R}$ by
  \[
    h(2k) := \bigl(\gamma \mapsto [[\gamma]](\omega_k)\bigr)
    ,\qquad
    h(2j+1) := g_j
    \qquad (j,k \in \mathbb{N})
    .
  \]
  Every $h(n)$ is measurable and $\overline\eta_l$-integrable: for even $n=2k$ this is Step~1, and
  for odd $n=2j+1$ this is the hypothesis on $g_j$.

  \emph{Step 3: the normalized probability measure.} Let $\mu := M^{-1} \bullet \overline\eta_l$
  (the measure $\overline\eta_l$ scaled by the extended-real constant $M^{-1}$). Since
  $\mu(\Theta(E)) = M^{-1}\cdot\overline\eta_l(\Theta(E)) = M^{-1}\cdot M = 1$ (using $M \ne 0,\infty$
  so that $M^{-1}\cdot M = 1$), $\mu$ is a probability measure on $\Theta(E)$. Scaling a measure by
  a finite nonzero extended-real constant preserves both measurability and integrability of a
  function, and rescales the integral by the same factor: each $h(n)$ ($n \in \mathbb{N}$) is
  therefore $\mu$-measurable and $\mu$-integrable, with
  \[
    \int_{\Theta(E)} h(n) \, d\mu = M_{\mathbb R}^{-1} \int_{\Theta(E)} h(n) \, d\overline\eta_l
    \qquad (n \in \mathbb{N})
    . \tag{$\ast$}
  \]

  \emph{Step 4: a full-$\mu$-measure subset of $\Theta_l(E) \cap \{\length=l\}$.} By
  \noderef{PSFamily} (field \emph{supported}) applied at this $l$, $\overline\eta_l(\Theta_l(E)^c)
  = 0$. By Mathlib's \texttt{exists\_measurable\_superset\_of\_null}, there is a measurable set
  $U \supseteq \Theta_l(E)^c$ with $\overline\eta_l(U) = 0$. Let
  $S := U^c \cap \{\gamma : \length(\gamma) = l\}$; $S$ is measurable, since $U$ is measurable and
  $\{\gamma : \length(\gamma)=l\}$ is measurable by \noderef{CurveLenMeasurable}. Since
  $\overline\eta_l(U) = 0$ and, by \noderef{PSFamily} (field \emph{lengthExactly}),
  $\overline\eta_l(\{\gamma : \length(\gamma) \ne l\}) = 0$, the union bound gives
  $\overline\eta_l(S^c) = \overline\eta_l\bigl(U \cup \{\length(\cdot)\ne l\}\bigr) \le
  \overline\eta_l(U) + \overline\eta_l(\{\length(\cdot)\ne l\}) = 0$, so $\overline\eta_l(S^c)=0$.
  Consequently $\mu(S^c) = M^{-1}\cdot\overline\eta_l(S^c) = M^{-1}\cdot 0 = 0$ as well.

  \emph{Step 5: apply \noderef{SamplingRealizesAverages}.} By Step~2 the family $(h(n))_{n \in
  \mathbb{N}}$ is measurable and $\mu$-integrable (Step~3), and by Step~4, $S$ is a measurable set
  with $\mu(S^c)=0$; \noderef{SamplingRealizesAverages} (with $X = \Theta(E)$, the probability
  measure $\mu$, the family $h$, and the set $S$) produces a sequence $(\gamma_i)_{i\in\mathbb{N}}$
  with $\gamma_i \in S$ for every $i$, and, for every $n \in \mathbb{N}$,
  \[
    \frac{1}{n}\sum_{i=0}^{n-1} h(n')(\gamma_i) \;\longrightarrow\; \int_{\Theta(E)} h(n')\,d\mu
    \qquad (n\to\infty)
    ,
    \tag{$\ast\ast$}
  \]
  simultaneously for every index $n' \in \mathbb{N}$ (writing $n'$ for the family index to avoid a
  clash with the averaging variable $n$).

  \emph{Conclusion of (i)--(ii).} Fix $i$. Since $\gamma_i \in S \subseteq \{\length(\cdot)=l\}$,
  $\length(\gamma_i) = l$, which is~(i). Since $\gamma_i \in S \subseteq U^c$ and
  $\Theta_l(E)^c \subseteq U$, $\gamma_i \notin \Theta_l(E)^c$, i.e.\ $\gamma_i \in \Theta_l(E)$,
  i.e.\ $\mathbb{M}([[\gamma_i]]) = l$ (\noderef{ThetaL}), which is~(ii).

  \emph{Conclusion of (iii).} Fix $k \in \mathbb{N}$ and apply $(\ast\ast)$ at $n' = 2k$:
  \[
    \frac{1}{n}\sum_{i=0}^{n-1} h(2k)(\gamma_i)
    = \frac{1}{n}\sum_{i=0}^{n-1} [[\gamma_i]](\omega_k)
    \;\longrightarrow\; \int_{\Theta(E)} h(2k)\,d\mu = \int_{\Theta(E)} [[\cdot]](\omega_k)\,d\mu
    .
  \]
  By $(\ast)$ at $n=2k$ and \noderef{PSFamily} (field \emph{weakLength}, which states
  $T(\omega_k) = l^{-1}\int_{\Theta(E)} [[\gamma]](\omega_k)\,d\overline\eta_l(\gamma)$, i.e.
  $\int_{\Theta(E)} [[\cdot]](\omega_k)\,d\overline\eta_l = l \cdot T(\omega_k)$),
  \[
    \int_{\Theta(E)} [[\cdot]](\omega_k)\,d\mu
    = M_{\mathbb R}^{-1} \cdot \int_{\Theta(E)} [[\cdot]](\omega_k)\,d\overline\eta_l
    = M_{\mathbb R}^{-1} \cdot l \cdot T(\omega_k)
    .
  \]
  Multiplying the displayed convergence by the constant $M_{\mathbb R}/l$ (which is legitimate
  since $l\ne0$, and preserves the limit by continuity of multiplication by a constant) gives
  \[
    \frac{M_{\mathbb R}}{n\cdot l}\sum_{i=0}^{n-1} [[\gamma_i]](\omega_k)
    \;\longrightarrow\;
    \frac{M_{\mathbb R}}{l}\cdot M_{\mathbb R}^{-1}\cdot l \cdot T(\omega_k)
    = T(\omega_k)
    ,
  \]
  using $M_{\mathbb R} \ne 0$ and $l \ne 0$ to cancel. This is exactly~(iii).

  \emph{Conclusion of (iv).} Fix $j \in \mathbb{N}$ and apply $(\ast\ast)$ at $n' = 2j+1$:
  \[
    \frac{1}{n}\sum_{i=0}^{n-1} h(2j+1)(\gamma_i) = \frac{1}{n}\sum_{i=0}^{n-1} g_j(\gamma_i)
    \;\longrightarrow\; \int_{\Theta(E)} h(2j+1)\,d\mu = \int_{\Theta(E)} g_j\,d\mu
    = M_{\mathbb R}^{-1}\int_{\Theta(E)} g_j\,d\overline\eta_l
    ,
  \]
  the last equality by $(\ast)$ at $n=2j+1$. This is exactly~(iv), completing the proof.
\end{proof}
